% Part:sets-functions-relations
% Chapter: sets
% Section: reduction

\documentclass[../../../include/open-logic-section]{subfiles}

\begin{document}

\olfileid[ar]{sfr}{siz}{red}

\olsection{الاختزال}

\begin{editorial}
نعيد هنا إثبات عدم القابلية للتعداد بطريق الاختزال، فتطابق النتائج ما في \olref[nen]{sec}. ومن أراد عرضًا أوجز قليلًا، موافقًا لنتائج \olref[nen-alt]{sec}، وجده في \olref[red-alt]{sec}.
\end{editorial}

قد ثبت بالحجة القطرية أن $\Pow{\PosInt}$ !!{nonenumerable}، وكان قد سبق ثبوت ذلك لـ$\Bin^\omega$، أي مجموعة المتتاليات اللانهائية من~${\OLMathNumeral{0}}$ و~${\OLMathNumeral{1}}$، فهي !!{nonenumerable}. ويمكن أن نجعل النتيجة الثانية سبيلًا آخر إلى إثبات أن $\Pow{\PosInt}$ !!{nonenumerable}: فنبرهن على أنه \emph{لو كانت $\Pow{\PosInt}$ !!{enumerable} لكانت $\Bin^\omega$ !!{enumerable} أيضًا}. وحيث علمنا أن $\Bin^\omega$ ليست !!{enumerable}، امتنع ذلك على $\Pow{\PosInt}$. وهذا هو \emph{اختزال} مسألة إلى أخرى؛ فإن مسألة تعداد $\Bin^\omega$ قد رددناها إلى مسألة تعداد $\Pow{\PosInt}$. فلو ظفرنا بحل الثانية، أعني تعداد $\Pow{\PosInt}$، لظفرنا بذلك بحل الأولى، أعني تعداد $\Bin^\omega$.

ومعنى اختزال تعداد مجموعة~${\OLId{latin-looped}{B}}$ إلى تعداد مجموعة~${\OLId{latin-looped}{A}}$ أن نجد طريقًا يحول تعداد~${\OLId{latin-looped}{A}}$ إلى تعداد~${\OLId{latin-looped}{B}}$. وأيسر ذلك أن نعين دالة~!!a{surjective} ${\OLId{latin-ordinary}{f}}\colon {\OLId{latin-looped}{A}} \to {\OLId{latin-looped}{B}}$. فمتى كانت ${\OLId{latin-ordinary}{x}}_{\OLMathNumeral{1}}$، ${\OLId{latin-ordinary}{x}}_{\OLMathNumeral{2}}$، \dots{} تعدادًا لـ~${\OLId{latin-looped}{A}}$، كانت ${\OLId{latin-ordinary}{f}}({\OLId{latin-ordinary}{x}}_{\OLMathNumeral{1}})$، ${\OLId{latin-ordinary}{f}}({\OLId{latin-ordinary}{x}}_{\OLMathNumeral{2}})$، \dots{} تعدادًا لـ~${\OLId{latin-looped}{B}}$. فمطلبنا في المسألة الحاضرة دالة شاملة ${\OLId{latin-ordinary}{f}}\colon \Pow{\PosInt} \to \Bin^\omega$.

\begin{prob}
لتكن دالة~!!{injective} ${\OLId{latin-ordinary}{g}}\colon {\OLId{latin-looped}{B}} \to {\OLId{latin-looped}{A}}$ موجودة، ولتكن ${\OLId{latin-looped}{B}}$~!!{nonenumerable}. بين أن ${\OLId{latin-looped}{A}}$ كذلك، بأن تبين كيف تحول~${\OLId{latin-ordinary}{g}}$ تعداد~${\OLId{latin-looped}{A}}$، لو وجد، إلى تعداد~${\OLId{latin-looped}{B}}$.
\end{prob}

\begin{proof}[برهان {\olref[nen]{thm:nonenum-pownat}} بالاختزال]
لنفرض أن $\Pow{\PosInt}$ !!{enumerable}، ونأخذ لها تعدادًا ${\OLId{latin-looped}{Z}}_{{\OLMathNumeral{1}}}$، ${\OLId{latin-looped}{Z}}_{{\OLMathNumeral{2}}}$، ${\OLId{latin-looped}{Z}}_{{\OLMathNumeral{3}}}$، \dots

ونضع الدالة ${\OLId{latin-ordinary}{f}} \colon \Pow{\PosInt} \to \Bin^\omega$ على الوجه الآتي: نجعل ${\OLId{latin-ordinary}{f}}({\OLId{latin-looped}{Z}})$ هي المتتالية ${\OLId{latin-ordinary}{s}}$ التي تكون فيها ${\OLId{latin-ordinary}{s}}({\OLId{latin-ordinary}{n}}) = {\OLMathNumeral{1}}$ إذا وفقط إذا كان ${\OLId{latin-ordinary}{n}} \in {\OLId{latin-looped}{Z}}$، وتكون ${\OLId{latin-ordinary}{s}}({\OLId{latin-ordinary}{n}}) = {\OLMathNumeral{0}}$ فيما عدا ذلك. وهذا تعيين لا التباس فيه؛ فمتى كان ${\OLId{latin-looped}{Z}} \subseteq \PosInt$، فإن ${\OLId{latin-ordinary}{n}} \in \PosInt$ إما أن يكون !!a{element} من ${\OLId{latin-looped}{Z}}$ وإما ألا يكون. فمجموعة الأعداد الصحيحة الموجبة الزوجية ${\OLMathNumeral{2}}\PosInt = \{{\OLMathNumeral{2}},
{\OLMathNumeral{4}}, {\OLMathNumeral{6}}, \dots\}$ صورتها ${\OLMathNumeral{010101}}\dots$، وصورة المجموعة الخالية ${\OLMathNumeral{0000}}\dots$، وصورة $\PosInt$ نفسها ${\OLMathNumeral{1111}}\dots$.

وهذه الدالة !!{surjective} أيضًا. ذلك أن كل متتالية من~${\OLMathNumeral{0}}$ و~${\OLMathNumeral{1}}$ تعين مجموعة الأعداد الصحيحة الموجبة التي يقع~${\OLMathNumeral{1}}$ في مواضعها، فتكون المتتالية صورة تلك المجموعة. ولتفصيله، خذ ${\OLId{latin-ordinary}{s}} \in \Bin^\omega$، وعين ${\OLId{latin-looped}{Z}}
\subseteq \PosInt$ بوضع
\[
{\OLId{latin-looped}{Z}} = \Setabs{{\OLId{latin-ordinary}{n}} \in \PosInt}{{\OLId{latin-ordinary}{s}}({\OLId{latin-ordinary}{n}}) = {\OLMathNumeral{1}}}
\]
فيكون ${\OLId{latin-ordinary}{f}}({\OLId{latin-looped}{Z}}) = {\OLId{latin-ordinary}{s}}$، كما يظهر بمراجعة تعريف~${\OLId{latin-ordinary}{f}}$.

فانظر الآن إلى القائمة
\[
{\OLId{latin-ordinary}{f}}({\OLId{latin-looped}{Z}}_{\OLMathNumeral{1}}), {\OLId{latin-ordinary}{f}}({\OLId{latin-looped}{Z}}_{\OLMathNumeral{2}}), {\OLId{latin-ordinary}{f}}({\OLId{latin-looped}{Z}}_{\OLMathNumeral{3}}), \dots
\]
إن ${\OLId{latin-ordinary}{f}}$ !!{surjective}، فلا يخلو عنصر من $\Bin^\omega$ من أن يكون قيمة لـ~${\OLId{latin-ordinary}{f}}$ عند وسيط ما؛ وكل وسيط وارد في التعداد المفروض، فلا بد من ورود تلك القيمة في هذه القائمة. فهي إذن تعداد يستوفي~$\Bin^\omega$.

وحاصل ذلك أنه لو كانت $\Pow{\PosInt}$ !!{enumerable} لكانت $\Bin^\omega$ !!{enumerable}. وقد ثبت أن $\Bin^\omega$ !!{nonenumerable} (\olref[nen]{thm:nonenum-bin-omega})، فلزم أن $\Pow{\PosInt}$ !!{nonenumerable}.
\end{proof}

\begin{editorial}
\textbf{تنبيه تصحيحي على الأصل (OLSIZ-004).}
سمّى الأصل المتتالية ${\OLId{latin-ordinary}{s}}_{\OLId{latin-ordinary}{k}}$ من غير أن يربط ${\OLId{latin-ordinary}{k}}$؛ فاستُعمل الاسم
${\OLId{latin-ordinary}{s}}$ في التعريف كله، إذ تحدد المجموعة ${\OLId{latin-looped}{Z}}$ متتاليتها المميزة وحدها.
\end{editorial}

\begin{explain}
وينبغي ضبط جهة الاختزال، إذ يقع اللبس فيها بسهولة. فوجود دالة~!!a{surjective} ${\OLId{latin-ordinary}{g}} \colon \Bin^\omega \to {\OLId{latin-looped}{B}}$ \emph{لا} يثبت أن ${\OLId{latin-looped}{B}}$ !!{nonenumerable}. وشاهده الدالة ${\OLId{latin-ordinary}{g}}
\colon \Bin^\omega \to \Bin$ التي نعرفها بـ${\OLId{latin-ordinary}{g}}({\OLId{latin-ordinary}{s}}) = {\OLId{latin-ordinary}{s}}({\OLMathNumeral{1}})$: فهي تأخذ المتتالية من~${\OLMathNumeral{0}}$ و~${\OLMathNumeral{1}}$ فترسلها إلى أول !!{element} فيها. وهي !!{surjective}، لأن متتاليات منها تبدأ بـ${\OLMathNumeral{0}}$ ومنها ما يبدأ بـ${\OLMathNumeral{1}}$؛ على أن $\Bin$ منتهية. كذلك لا بد من كون الدالة~${\OLId{latin-ordinary}{f}}$ !!{surjective}، وإلا لم يتم البرهان؛ إذ لا يلزم حينئذ أن تستوفي ${\OLId{latin-ordinary}{f}}({\OLId{latin-ordinary}{x}}_{\OLMathNumeral{1}})$، ${\OLId{latin-ordinary}{f}}({\OLId{latin-ordinary}{x}}_{\OLMathNumeral{2}})$، \dots{} !!{element}s~${\OLId{latin-looped}{B}}$ كلها. وانظر مثلًا إلى العبارة
\[
{\OLId{latin-ordinary}{h}}({\OLId{latin-ordinary}{n}}) = \underbrace{{\OLMathNumeral{00}}\dots{\OLMathNumeral{0}}}_{\text{${\OLId{latin-ordinary}{n}}$ أصفار}}{\OLMathNumeral{111}}\dots
\]
أي إن ${\OLId{latin-ordinary}{h}}({\OLId{latin-ordinary}{n}})$ متتالية لا نهائية تبدأ بـ${\OLId{latin-ordinary}{n}}$ أصفار ثم تستمر بالآحاد.
وهذا يعرّف دالة ${\OLId{latin-ordinary}{h}}\colon \PosInt \to \Bin^\omega$، مع أن
$\PosInt$ !!{enumerable}.
\end{explain}

\begin{editorial}
\textbf{تنبيه تصحيحي على الأصل (OLSIZ-005).}
جعل الأصل قيمة ${\OLId{latin-ordinary}{h}}({\OLId{latin-ordinary}{n}})$ سلسلة منتهية من الأصفار مع أن المجال المقابل
$\Bin^\omega$ لا يضم إلا المتتاليات اللانهائية؛ فأضيف ذيل لا نهائي من الآحاد.
\end{editorial}

\begin{prob}
أقم حجة اختزال على أن مجموعة جميع \emph{المجموعات المؤلفة من} أزواج الأعداد الصحيحة الموجبة !!{nonenumerable}.
\end{prob}

\begin{prob}\label{sfr:siz:red:prob:nat-nat}
بين بالاختزال أن المجموعة~${\OLId{latin-looped}{X}}$ التي تضم جميع الدوال ${\OLId{latin-ordinary}{f}}\colon \Nat \to \Nat$ !!{nonenumerable}. وتلميحه أن تعين دالة شاملة من ${\OLId{latin-looped}{X}}$ إلى~$\Bin^\omega$.
\end{prob}

\begin{prob}
استعمل الاختزال لإثبات أن $\Nat^\omega$ !!{nonenumerable}؛ والمراد بها مجموعة جميع المتتاليات اللانهائية من الأعداد الطبيعية.
\end{prob}

\begin{prob}
سمّ ${\OLId{latin-looped}{P}}$ مجموعة الدوال من الأعداد الصحيحة الموجبة إلى $\{{\OLMathNumeral{0}}\}$، وسمّ ${\OLId{latin-looped}{Q}}$ مجموعة الدوال \emph{الجزئية} من الأعداد الصحيحة الموجبة إلى $\{{\OLMathNumeral{0}}\}$. والمطلوب إثبات أن ${\OLId{latin-looped}{P}}$~!!{enumerable}، بخلاف ${\OLId{latin-looped}{Q}}$. وتلميحه أن ترد مسألة تعداد $\Bin^\omega$ إلى تعداد~${\OLId{latin-looped}{Q}}$.
\end{prob}

\begin{prob}
إذا كانت ${\OLId{latin-looped}{S}}$ مجموعة الدوال~!!{surjective} من الأعداد الصحيحة الموجبة إلى \{\OLClassicalDigits{0},\OLClassicalDigits{1}\}، أي كانت ${\OLId{latin-looped}{S}}$ تضم جميع الدوال~!!{surjective}~${\OLId{latin-ordinary}{f}}\colon \PosInt \to \Bin$، فأثبت أن ${\OLId{latin-looped}{S}}$ !!{nonenumerable}.
\end{prob}

\begin{prob}
بين أن $\Real$، مجموعة الأعداد الحقيقية بأسرها، !!{nonenumerable}.
\end{prob}

\end{document}
